% =============================================================================
%  WorkSheet Generator (wsg) -- default macros
%
%  Mathematical notation shared by the worksheets. The set is taken from the
%  lecture notes of Applied Mathematics (github.com/D4vEOFF/AM-skripta) so
%  that a task can be copied from the notes into a worksheet unchanged.
%
%  A copy of this file is placed into every new project (wsg new) and is
%  pulled in by the header template through the <<MACROS>> placeholder --
%  after the packages and before the worksheet machinery. Nothing that belongs
%  to the look of the sheet (task, solution, page style) is defined here.
% =============================================================================

% -----------------------------------------------------------------------------
%  Quotation marks of the current language: \enquote{...}
%
%  csquotes can also turn " into an active character which opens and closes
%  the quotation automatically. It is left switched off on purpose -- an active
%  quote clashes with verbatim material and breaks inside the discarded
%  solutions of the assignment sheet. Enable it per project if you want it:
%
%    \MakeOuterQuote{"}
% -----------------------------------------------------------------------------

% -----------------------------------------------------------------------------
%  Number sets and basic notation
% -----------------------------------------------------------------------------
\DeclareDocumentCommand{\R}{}{\mathbb{R}}
\DeclareDocumentCommand{\C}{}{\mathbb{C}}
\DeclareDocumentCommand{\N}{}{\mathbb{N}}
\DeclareDocumentCommand{\Q}{}{\mathbb{Q}}
\DeclareDocumentCommand{\Z}{}{\mathbb{Z}}
\DeclareDocumentCommand{\F}{}{\mathbb{F}}

\renewcommand{\emptyset}{\varnothing}

\DeclareDocumentCommand{\set}{m}{\left\{#1\right\}}
\DeclareDocumentCommand{\powset}{m}{\mathcal{P}(#1)}
\DeclareDocumentCommand{\sizeof}{m}{\left| #1 \right|}
\DeclareDocumentCommand{\abs}{m}{\left|\,#1\,\right|}
\DeclareDocumentCommand{\setcomplement}{m}{\ensuremath{{#1}^\mathrm{c}}}
% Separator inside a set: \set{x\in\R\admid x>0}
\DeclareDocumentCommand{\admid}{}{\mathrel{:}}
% Solution set of a system of equations
\DeclareDocumentCommand{\solutions}{m}{\left[K=\left\{#1\right\}\right]}

\newcommand{\map}[3]{\ensuremath{#1:#2\rightarrow #3}}
\newcommand{\floor}[1]{\left\lfloor#1\right\rfloor}
\newcommand{\ceil}[1]{\left\lceil#1\right\rceil}
\newcommand{\eps}{\varepsilon}
\newcommand{\hmid}{\mathrel{\rule[0.6ex]{1.2em}{0.4pt}}}
\DeclareDocumentCommand{\degre}{}{\ensuremath{^\circ}}

% -----------------------------------------------------------------------------
%  Operators
% -----------------------------------------------------------------------------
\DeclareMathOperator{\tg}{tg}
\DeclareMathOperator{\cotg}{cotg}
\DeclareMathOperator{\arctg}{arctg}
\DeclareMathOperator{\arccotg}{arccotg}
\DeclareMathOperator{\sgn}{sgn}
\DeclareMathOperator{\character}{char}
\DeclareMathOperator{\id}{id}
\DeclareMathOperator{\dom}{dom}
\DeclareMathOperator{\image}{Im}
\DeclareMathOperator{\kernel}{Ker}
\DeclareMathOperator{\Hom}{Hom}
\DeclareMathOperator*{\argmin}{arg\,min}
\DeclareMathOperator*{\argmax}{arg\,max}

\newcommand{\bigO}{\mathcal{O}}
\newcommand{\bigOmega}{\Omega}

% Logic gates
\DeclareMathOperator{\NOT}{NOT}
\DeclareMathOperator{\AND}{AND}
\DeclareMathOperator{\NAND}{NAND}
\DeclareMathOperator{\OR}{OR}
\DeclareMathOperator{\NOR}{NOR}
\DeclareMathOperator{\XOR}{XOR}

% -----------------------------------------------------------------------------
%  Linear algebra
% -----------------------------------------------------------------------------
\DeclareMathOperator{\Dim}{dim}
\DeclareMathOperator{\Span}{span}
\DeclareMathOperator{\Kan}{kan}
\DeclareMathOperator{\REF}{REF}
\DeclareMathOperator{\RREF}{RREF}
\DeclareMathOperator{\rank}{rank}
\DeclareMathOperator{\rowspace}{\mathcal{R}}
\DeclareMathOperator{\colspace}{\mathcal{S}}

\DeclareDocumentCommand{\matrow}{ m m }{\ensuremath{{#1}_{#2*}}}
\DeclareDocumentCommand{\matcol}{ m m }{\ensuremath{{#1}_{*#2}}}
\DeclareDocumentCommand{\transpose}{ m }{\ensuremath{{#1^{\top}}}}

% Coordinates with respect to a basis and the matrix of a homomorphism
\DeclareDocumentCommand{\coord}{ m m }{\ensuremath{\left[#1\right]_{#2}}}
\DeclareDocumentCommand{\hommat}{ m m m }{\ensuremath{{}_{#1}\!\left[#2\right]_{#3}}}

% -----------------------------------------------------------------------------
%  Error correcting codes
% -----------------------------------------------------------------------------
\DeclareDocumentCommand{\hweight}{m}{\ensuremath{w\!\left(#1\right)}}
\DeclareDocumentCommand{\hdist}{mm}{\ensuremath{d\!\left(#1,#2\right)}}
\DeclareDocumentCommand{\mindist}{m}{\ensuremath{d\!\left(#1\right)}}

% -----------------------------------------------------------------------------
%  Marks and highlighting
% -----------------------------------------------------------------------------
\DeclareDocumentCommand{\cmark}{}{\ding{51}}
\DeclareDocumentCommand{\xmark}{}{\ding{55}}
\NewDocumentCommand{\markred}{m}{\textcolor{red}{#1}}
\NewDocumentCommand{\markblue}{m}{\textcolor{blue}{#1}}
\NewDocumentCommand{\todo}{m}{\textcolor{red}{\textbf{TODO:} #1}}
\DeclareDocumentCommand{\rightnote}{m}{\hspace*{\fill} $\triangleleft$ \textit{#1}}

% A circled symbol that keeps the size of its content
\makeatletter
\newcommand{\circled}[1]{\mathpalette\circled@{#1}}
\newcommand{\circled@}[2]{%
  \tikz[baseline=(X.base)]{%
    \node[inner sep=0pt, outer sep=0pt] (X) {$#1#2$};
    \begin{scope}[overlay]
      \node[draw, circle, fit=(X), inner sep=0.35ex] {};
    \end{scope}
  }%
}
\makeatother

% -----------------------------------------------------------------------------
%  Statement of a computational problem
% -----------------------------------------------------------------------------
\DeclareDocumentCommand{\problem}{mmm}{%
  \begingroup
  \setlength{\fboxsep}{3mm}%
  \fbox{%
    \begin{minipage}{0.65\linewidth}
      \textsc{#1}\\[3mm]
      \indent\textbf{Vstup:} #2\\[3mm]
      \indent\textbf{Výstup:} #3
    \end{minipage}%
  }%
  \endgroup
}

% -----------------------------------------------------------------------------
%  Statements -- useful when a definition has to be recalled on the sheet
% -----------------------------------------------------------------------------
\theoremstyle{definition}
\newtheorem{definition}{Definice}
\newtheorem*{definition*}{Definice}
\newtheorem{theorem}{Věta}
\newtheorem*{theorem*}{Věta}
\newtheorem{lemma}[theorem]{Lemma}
\newtheorem{proposition}[theorem]{Tvrzení}
\newtheorem{corollary}[theorem]{Důsledek}
\newtheorem{example}{Příklad}
\newtheorem*{example*}{Příklad}
\newtheorem{remark}{Poznámka}
\newtheorem*{remark*}{Poznámka}
\newtheorem*{symboldef}{Značení}
\AtBeginEnvironment{remark}{\footnotesize}

% -----------------------------------------------------------------------------
%  TikZ: vertices and edges of a graph
% -----------------------------------------------------------------------------
\tikzset{
  vertex/.style={thick,draw,circle,minimum size=4mm,inner sep=1pt,
    outer sep=1pt,align=center},
  vertexplain/.style={minimum size=4mm,inner sep=1pt,outer sep=1pt,align=center},
  verteximg/.style={draw,circle,minimum size=4mm,inner sep=0pt,
    outer sep=1pt,align=center,clip},
  verteximgplain/.style={minimum size=4mm,inner sep=0pt,outer sep=1pt,align=center},
  edge/.style={thick,>=Stealth},
  diredge/.style={thick,->,>=Stealth}
}
